\section{Problem}

% ============================================================
% SLIDE 2
% ============================================================
\begin{frame}[shrink=11]{Một bảng dữ liệu, ba câu hỏi}{Problem \textperiodcentered\ motivation}
\begin{columns}[T,onlytextwidth]
\begin{column}{0.43\textwidth}
  {\large\bfseries\color{DTNavy}Ví dụ minh họa: lượt tìm kiếm điện thoại theo tuần}
  \vspace{3mm}
  \begin{itemize}
    \item Hai từ khóa: \textbf{iPhone}, \textbf{Galaxy}.
    \item Ba khu vực: Bắc Mỹ, châu Âu, Đông Nam Á.
    \item Quan sát mới đến sau mỗi tuần.
  \end{itemize}
  \vspace{2mm}
  \begin{intuition}[Không chỉ là đếm]
    Ta muốn hiểu \textbf{quy luật đang thay đổi}, không chỉ nhìn một tổng cuối kỳ.
  \end{intuition}
\end{column}
\begin{column}{0.55\textwidth}
  \begin{questioncard}
    \softlabel{01 \quad TĂNG HAY GIẢM?}\par
    {\large Mức quan tâm đang đi theo hướng nào?}\par
    {\small\color{DTMuted}Trend khác nhau theo từ khóa và khu vực.}
  \end{questioncard}
  \vspace{2mm}
  \begin{questioncard}
    \softlabel{02 \quad CÓ LAN GIỮA CÁC VÙNG?}\par
    {\large Nơi nào đi trước, nơi nào thay đổi sau?}\par
    {\small\color{DTMuted}Ta cần diễn giải liên hệ, không chỉ fit từng chuỗi.}
  \end{questioncard}
  \vspace{2mm}
  \begin{questioncard}
    \softlabel{03 \quad TUẦN TỚI RA SAO?}\par
    {\large Có thể dự báo khi dữ liệu vẫn đang đến?}\par
    {\small\color{DTMuted}Mô hình cũ có thể lỗi thời ngay khi quy luật đổi.}
  \end{questioncard}
\end{column}
\end{columns}
\vspace{2mm}
\papersource{ý tưởng marketing trong §1, p. 1; chi tiết iPhone/Galaxy và ba vùng là ví dụ giảng dạy.}
\end{frame}

% ============================================================
% SLIDE 3
% ============================================================
\begin{frame}[shrink=17]{Dữ liệu thực chất là một chồng bảng theo tuần}{Problem \textperiodcentered\ tensor bằng ví dụ}
\begin{columns}[c,onlytextwidth]
\begin{column}{0.50\textwidth}
\centering
\begin{tikzpicture}[x=.78cm,y=.68cm]
  \foreach \shift/\week/\shade in {0/t/35,1.0/{t+1}/60} {
    \begin{scope}[shift={(\shift,-\shift)}]
      \fill[white,draw=DTLine] (-.12,-.12) rectangle (3.45,2.55);
      \foreach \x/\v in {0/25,1/50,2/75} {
        \fill[DTTeal!\v,draw=white] (\x,1.15) rectangle ++(1,1);
        \fill[DTTeal!\shade,draw=white] (\x,0) rectangle ++(1,1);
      }
      \node[tinylabel,anchor=west] at (3.55,1.05) {tuần $\week$};
    \end{scope}}
  \node[tinylabel,anchor=east] at (-.48,1.68) {iPhone};
  \node[tinylabel,anchor=east] at (-.48,.52) {Galaxy};
  \node[tinylabel] at (.52,-1.45) {Bắc Mỹ};
  \node[tinylabel] at (1.52,-1.45) {châu Âu};
  \node[tinylabel] at (2.52,-1.45) {ĐNÁ};
  \draw[dtarrow] (4.7,-1.45) -- (5.55,-1.45);
  \node[tinylabel,anchor=west,text=DTGold] at (5.65,-1.45) {thời gian};
\end{tikzpicture}
\vspace{3mm}

{\small\color{DTMuted}Hai lát 2 × 3; màu đậm hơn biểu thị giá trị lớn hơn (minh họa).}
\end{column}
\begin{column}{0.47\textwidth}
  \begin{intuition}[Nói đơn giản]
    Mỗi tuần là một \textbf{bảng từ khóa × địa điểm}. Xếp các bảng theo thứ tự thời gian, ta có \textbf{tensor bậc ba}.
  \end{intuition}
  \vspace{3mm}
  \begin{mathcard}[Ký hiệu]
    \centering
    $X_t\in\mathbb N^{k\times l}$ \quad(một tuần)\\[2mm]
    $\mathcal X\in\mathbb N^{k\times l\times t_c}$ \quad(đến hiện tại)
  \end{mathcard}
  \vspace{3mm}
  \begin{itemize}
    \item Ví dụ: $k=2$ từ khóa, $l=3$ vùng.
    \item Một tuần có $2\times3=6$ ô.
    \item 104 tuần có $2\times3\times104=624$ ô minh họa.
  \end{itemize}
\end{column}
\end{columns}
\vspace{1mm}
\papersource{Problem definition, §3.1, p. 4; kích thước 2 × 3 × 104 là minh họa.}
\end{frame}

% ============================================================
% SLIDE 4
% ============================================================
\begin{frame}[shrink=24]{Problem 1 --- stream forecasting}{Problem \textperiodcentered\ input và output}
\centering
\begin{tikzpicture}[x=.83cm,y=.75cm]
  \draw[DTLine,line width=2pt] (0,0) -- (13,0);
  \foreach \x/\lab in {0/$t_1$,1.4/$t_2$,2.8/$\cdots$,8.6/$t_{c-1}$,10/$t_c$,12.6/$t_c+L_f$} {
    \fill[DTTeal] (\x,0) circle (.08);
    \node[tinylabel,anchor=north] at (\x,-.28) {\lab};}
  \draw[DTGold,line width=4pt] (4.3,.4) -- (10,.4);
  \draw[DTGold] (4.3,.1) -- (4.3,.72); \draw[DTGold] (10,.1) -- (10,.72);
  \node[biglabel,anchor=south] at (7.15,.72) {cửa sổ hiện tại $\mathcal X^c$: $L_c$ bước};
  \draw[dtarrow] (10.25,0) -- (12.25,0);
  \node[font=\sffamily\small,text=DTGold,anchor=south] at (11.2,.1) {dự báo};
\end{tikzpicture}
\vspace{3mm}

\begin{columns}[T,onlytextwidth]
\begin{column}{0.48\textwidth}
\begin{mathcard}[Input tại mỗi thờingian]
  \centering
  $\mathcal X^c=\{X_t\}_{t=t_c-L_c+1}^{t_c}$\\[2mm]
  {\small Cửa sổ gần nhất thay vì toàn bộ lịch sử.}
\end{mathcard}
\end{column}
\begin{column}{0.48\textwidth}
\begin{mathcard}[Output]
  \centering
  time-varying dynamics $+\widehat X_{t_c+L_f}$\\[2mm]
  {\small Trend, seasonality, diffusion và dự báo $L_f$ bước.}
\end{mathcard}
\end{column}
\end{columns}
\vspace{3mm}
\begin{takeaway}[Ví dụ cấu hình paper]
  GoogleTrends dùng $L_c=104$ tuần và các tầm dự báo 13/26/39 tuần; Covid-19 dùng $L_c=56$ ngày và 7/14/21 ngày.
\end{takeaway}
\vspace{2mm}
\papersource{Problem 1, §3.1, p. 4; experimental protocol, §5 Q1, p. 7.}
\end{frame}

% ============================================================
% SLIDE 5
% ============================================================
\begin{frame}[shrink=25]{Vì sao khó? Ba phức tạp chồng lên nhau}{Problem \textperiodcentered\ ba thách thức}
\begin{columns}[T,onlytextwidth]
\begin{column}{0.32\textwidth}
\begin{intuition}[01 \quad Complex patterns]
  \textbf{Trend} tăng/giảm, \textbf{seasonality} lặp, \textbf{diffusion} lan và \textbf{outlier} bất ngờ cùng tồn tại.\\[2mm]
  {\small Cường độ và thờingian khác theo keyword/location.}
\end{intuition}
\end{column}
\begin{column}{0.32\textwidth}
\begin{mathcard}[02 \quad Dynamic changes]
  Quy luật \textbf{thay đổi theo thời gian}.\\[2mm]
  Một diffusion pattern có thể chỉ xuất hiện trong giai đoạn ngắn; model học từ quá khứ có thể lỗi thời ngay.
\end{mathcard}
\end{column}
\begin{column}{0.32\textwidth}
\begin{caution}[03 \quad High-dimensional stream]
  Nhiều từ khóa × nhiều quốc gia × thời gian không giới hạn.\\[2mm]
  Tăng chi phí tính toán và làm kết quả khó diễn giải.
\end{caution}
\end{column}
\end{columns}
\vspace{5mm}
\centering
\begin{tikzpicture}[x=.85cm,y=.75cm]
  \node[dtbox,minimum width=2.8cm,minimum height=.9cm] (a) at (0,0){nhiều chuỗi};
  \node[dtbox,minimum width=2.8cm,minimum height=.9cm] (b) at (4,0){nhiều quy luật};
  \node[dtbox,minimum width=2.8cm,minimum height=.9cm] (c) at (8,0){nhiều giai đoạn};
  \draw[dtarrow] (a)--(b); \draw[dtarrow] (b)--(c);
  \node[dtgold,minimum width=4.2cm,minimum height=1.0cm,align=center] (d) at (4,-2)
    {cần \textbf{nén + tách + thích nghi}};
  \draw[dtdash] (a.south)--(d.west); \draw[dtdash] (c.south)--(d.east);
\end{tikzpicture}
\vspace{2mm}
\papersource{§1, pp. 1--2.}
\end{frame}

% ============================================================
% SLIDE 6
% ============================================================
\begin{frame}[shrink=5]{Lộ trình buổi trình bày}{Problem \textperiodcentered\ roadmap}
\begin{columns}[T,onlytextwidth]
\begin{column}{0.235\textwidth}
\begin{intuition}[01 \quad Method]
  Tách và nén dữ liệu; mô hình hóa reaction--diffusion; học và cập nhật.
\end{intuition}
\end{column}
\begin{column}{0.235\textwidth}
\begin{mathcard}[02 \quad Paper experiments]
  Accuracy, scalability, effectiveness trên sáu tensor streams.
\end{mathcard}
\end{column}
\begin{column}{0.235\textwidth}
\begin{questioncard}
  \softlabel{03 \quad REPRODUCTION}\par
  Protocol, bảng Paper vs. nhóm, runtime và lessons.\\[2mm]
  {\small Phần kết quả chờ thực nghiệm của nhóm.}
\end{questioncard}
\end{column}
\begin{column}{0.235\textwidth}
\begin{takeaway}[04 \quad Demo]
  Một slide trống để nhóm thêm sản phẩm minh họa.
\end{takeaway}
\end{column}
\end{columns}
\vspace{7mm}
\centering
{\large\color{DTNavy}Xuyên suốt method: quay lại ba câu hỏi \textbf{tăng/giảm --- lan --- dự báo}.}
\end{frame}
